\chapter{42}



Louis hob den bepackten Rucksack vom Bett. Manuela hatte ihm einige
Dinge gereicht, die er auf der Alp benötigen würde. Die Liste mit
Notizen steckte er in eine der beiden Seitentaschen seiner Jacke. Manuela hatte ihm
geduldig seine vorläufigen Fragen beantwortet.

«Weitere Details können wir beim Abendessen klären. Joh wird dich heute
Abend vor allem in Sachen Berghütte instruieren, und der Rest ergibt
sich, wenn du oben bist. Und ja, Mazi hilft dir, wo immer nötig. Ihr
werdet euch bestimmt verstehen.»

Louis setzte sich aufs Bett und zog sein Handy aus der Hosentasche. Es
war soweit. Er stellte die Nummer seiner Mutter ein. Sein Herzschlag schnellte
hoch. Sein Hals wurde trocken. Wie das letzte Mal konkurrierten
Aufregung, Neugierde und Beruhigung um seine Zuwendung.
Die Verbindung klappte nicht. Sie war nicht erreichbar. Ihr Handy wohl
ausgeschaltet. Louis hielt sich die Hände vors Gesicht und rieb sich die 
Augen. Ach, Mensch, flüsterte er vor sich hin und presste die Lippen
zusammen. Er kippte zur Seite und hob die Beine aufs Bett. Würde es ihm gelingen, 
seine wichtigsten Fragen zu klären? Er musste unauffällig an Informationen zu Giorgias Zustand 
kommen und sicher sein, dass seine Mutter den Facebook-Aufruf gelöscht hatte.
Er wollte sie beruhigen und sie mit ihren Fragen auf später vertrösten. Wie
jedes Mal, wenn er an sie dachte, zog sich sein Herz 
zusammen. Es half, sich vorzustellen, wie er später schliesslich ihre Hand halten würde. Wie sie verstehen würde,
dass es unmöglich gewesen war, sie einzuweihen. Weil er sie doch nicht
in Schwierigkeiten hatte bringen wollen. Und dass er alles erstmal selbst
hatte klären müssen, bevor er ihr von dieser Katastrophe erzählen
konnte. Sie würde es verstehen. So war sie. Und dann, schwor er sich, musste
Schluss sein damit, ihr Sorgen zu bereiten.

Beim dritten Versuch nahm sie den Anruf entgegen.

«Pronto», ihre Stimme wackelte.

«Maman? Ich bin’s, Louis,» begann er und atmete ein, «wie geht’s dir?»

«Louis? Louis.»

Sie wirkte überrascht und aufgeregt.

«Ja, ich bin’s, Maman, mir geht’s gut. Ich brauch noch etwas Zeit und - »

«Ja, ich weiss, aber wo bist du denn, Louis? Ich versteh es noch immer
nicht, keiner hier versteht, was mit dir los ist.»

«Mir geht es gut. Sobald ich alles gelöst habe, werde ich dir’s
erzählen. Aber sag, das letzte Mal sprachst du von einem Unfall?»

Maman schien nach Luft zu schnappen. Sie sprach hektisch, als fehle ihr
die Zeit.

«Ja, Giorgias Mutter rief mich an. Es geht ihr schlecht. Giorgia ist
immer noch im Koma - ,» bei Giorgias Namen schluckte Louis. Maman klang besorgt. Sie hatten sich
über die Zeit ins Herz geschlossen, ihre beiden Mütter, « - die Salas sind voller Sorge, eine Rücken-OP steht an, sobald sie stabiler ist, vielleicht wird sie - also ich weiss nicht, wie schwer sie verletzt ist.»

Louis riss sich zusammen und gab den Ahnungslosen.

«Das ist ja schrecklich, Maman, die arme Giorgia, aber, was ist denn
überhaupt passiert?»

«Das weiss niemand genau, die Polizei hält sich bedeckt, habe ich
gehört, und Bastiano wurde vernommen. Valentina hat’s mir erzählt, sie
meinte, dass es Spuren gab an der Unfallstelle, von ihm.»

Louis liessen ihre Worte blitzartig in sich zusammenfallen. Seine Beine
wurden schwer.

«Ach ja?» warf er ein.

Maman schien mit sich zu ringen.

«Ja. Ich bin sehr traurig für Giorgia und auch für die Salas, einfach
furchtbar alles.»

Sie seufzte.

«Ja, und man sagt, dass Bastiano kein unbeschriebenes Blatt ist.»

Ihre Stimme wurde noch wackeliger. Louis glaubte, eine leise Wut
herauszuhören. Vielleicht war es seine eigene.

«Wie meinst du das?» fragte er vorsichtig.

«Der ist vorbestraft, also als Jugendlicher war er straffällig, darum
verschwand er ja in ein Internat. Wusstest du das? Und später,» Louis
hörte sie erneut nach Atem ringen, «später muss er wiederholt kriminell
geworden sein.»

«Nein, das wusste ich nicht.»

Er erstarrte. Dass Bastiano damals aufgrund krimineller Handlungen
verschwunden war, war ihm nicht bekannt. Ob Giorgia ihm etwas
verschwiegen, oder es auch selbst nicht gewusst hatte?

«Aber Louis, was ist mit dir? Wohin bist du denn verschwunden? Und
warum?»

Ihre Stimme hörte sich nun traurig an. Louis konnte sich gut vorstellen, wie
besorgt sie war. Und doch, er war ein erwachsener Mann. Darüber hatten
sie immer mal wieder gesprochen. Intensiv damals, als Louis bei Marco zu
arbeiten begonnen hatte.

«Wie ich dir sagte, Maman, ich muss was klären, ich brauche Zeit für
mich. Kein Grund zur Sorge. Ich mach das alles selbst - und Maman,»
dabei versuchte Louis in den Lautsprecher hinein zu lächeln, womit er sich
in erster Linie selbst beruhigte, «ich bin erwachsen, okay?»

«Ja, schon, Louis,» sie klang beschwichtigend, «aber du verstehst auch,
wie unerklärlich dein Verschwinden für mich ist? Und dass du mit einer
unterdrückten Nummer anrufst, wirkt, wie soll ich sagen, ungewöhnlich,
wir sind uns doch - .»

«Ja, das kann ich nachvollziehen. Aber den Facebook-Aufruf hast du
gelöscht, oder?»

«Ja, nur - .»

«Du wirst mich später verstehen, wenn ich dir alles erzähle. Aber Maman,
sag, wie geht es dir mit deinem \emph{Brando}?»

Sie fing überraschend an zu kichern.

«Jetzt hör mal auf, Louis,» sie machte eine kurze Pause und spielte die
Verlegene, «nenn ihn nicht dauernd \emph{Brando,»} sie zog den Namen in
die Länge, «Mauro meldet sich oft, wir gehen immer mal wieder aus, aber
das reicht im Moment.»

«Das müsste dann schon mal griffiger werden, oder?»

Louis musste über seine eigenen Worte lachen. Er hörte sie durchatmen.
Und freute sich. Gerade redeten sie zusammen wie immer. 

«Spass beiseite, Maman, ich wünsch dir echt Glück mit ihm.»

«Eins nach dem andern, schauen wir mal. Aber, Louis - ich kenn dich doch, du versuchst vom Wesentlichen abzulenken. Dass du immer mal wieder abtauchst, kenn ich ja, aber diesmal, diesmal mach ich mir wirklich Sorgen, weil - »

«Ja, hm, also, Maman, ich verstehe dich doch, aber lass mich bitte machen. Ich melde mich wieder. Und jetzt muss ich Schluss machen. Ich liebe dich.»

«Ich lieb dich auch, Louis. Und, und bitte melde dich bald wieder und pass gut auf dich auf, ja?»

«Klar, und du auf dich, Maman.»


\sectionbreak

Louis streckte die Arme von sich. Das Gespräch war nicht schlecht gelaufen. Zwischendurch hatte sie sogar gelacht.
Die Informationen zu Giorgias Zustand hingegen waren verstörend. Wenn
sie im Koma lag, dann war sie vielleicht beim Auffinden nicht
ansprechbar gewesen. Hatte Giorgia womöglich nichts über den Hergang
erzählen können. Sie sich in einem Krankenhausbett liegend vorzustellen, war unerträglich. 
Und was Maman von Bastiano wusste, verblüffte ihn. Wie konnte das überhaupt möglich sein? Bastiano
hatte sich doch bereits viel früher verabschiedet.
Mit keinem Wort hatte sie sein Verschwinden mit Giorgias Situation
in Verbindung gebracht. Nichts deutete darauf hin, dass man ihn suchte.
Oder verschwieg sie ihm etwas? Seine Mutter? Nein, dafür standen sie
sich zu nahe. Seine Gedanken begannen wieder unordentlich durcheinander
zu geraten.
Johs dunkle Stimme von unten war die willkommene Erlösung.

«Ciro, Einladung nach oben, Essen ist bereit.»

Schon alleine seine Stimme beruhigte Louis. Joh musste von seinem
Ausflug nach Brig zurück sein.

«Gerne,» rief Louis laut, «komme sofort.»

Damit schwang er seine Beine aus dem Bett, stellte sich vor den
Waschtisch, wusch sich die Hände und fuhr sich übers Gesicht.

«Dio mio,» entfuhr es ihm und Louis schaute sich im Spiegel in die Augen,
«was kommt als Nächstes?»
